% usobasico.tex
%
% Copyright (C) 2022--2026 José A. Navarro Ramón <janr.devel@gmail.com>
% Licencia del código GPLv2
% Licencia Creative Commons Recognition Non-Commercial Share-alike.
% (CC-BY-NC-SA)


\chapter{Utilizando {\sf{chemfig}}}
Este capítulo tiene los siguientes objetivos:
\begin{enumerate}
\item Aprender cómo cargar el paquete en Lua\LaTeX.
\item Familiarizarse con sus principales características, sobre todo las que se aplican
  en un contexto de ESO y bachillerato.
\item Conocer algunas formas de personalizar las fórmulas y/o reacciones químicas con
  respecto a las que presenta por defecto el paquete.
\item Conocer sus puntos fuertes y carencias.
\end{enumerate}

\section{Introducción}
En esta sección veremos algunos detalles importantes para usar este paquete.

\subsection{¿Cuándo interesa utilizarlo?}
Este paquete de química es más complejo que \textsf{mhchem}
Sirve para escribir:
\begin{enumerate}
\item Compuestos orgánicos (o inorgánicos, sobre todo cuando se muestran los enlaces).
\item Reacciones químicas complejas donde aparecen compuestos orgánicos e inorgánicos.
\end{enumerate}

Se recomienda utilizar \textsf{mhchem} en las situaciones que se mostraron en el
capítulo anterior. Por tanto, es recomendable repasar el paquete del capítulo anterior.
En algunos de esos casos también se podría utilizar también \textsf{chemfig}, aunque
teniendo que escribir más. En algunas de esas situaciones especiales podría interesar
\textsf{chemfig}, personalizando algo la fórmula.

En realidad, \textsf{chemfig} tiene muchas formas de personalizar las fórmulas y/o
reacciones químicas, mientras que \textsf{mhchem} es más limitado en ese sentido.




%%% Local Variables:
%%% coding: utf-8
%%% mode: latex
%%% TeX-engine: luatex
%%% TeX-master: "../chemfigtut.tex"
%%% End:
